\section{Conclusion and Roadmap}
\label{sec:conclusion}

The framework presented here is a small library of separately specified, separately
validated economic mappings between a least-cost energy--land--water plan and an
overlapping-generations economy. Each mapping carries one mechanism, enters the economy
on one account, and is tested at the layer of the claim it makes; the coupled
application composes them without hiding them, so any aggregate result can be traced
back to the mapping --- and the assumption --- that produced it. Relative to the
aggregate energy--economy couplings this literature descends from, the economy side
here prices the consequences onto households that differ by age and lifetime income,
onto industries that differ in exposure, and onto a government that must finance what
the plan implies --- which is what makes incidence, cohort, and fiscal questions
answerable at all.

The path from here is concrete, and each step removes a named limitation:
\begin{enumerate}
\item \textbf{A marginal price signal.} The commodity dual is degenerate in the
  present case, so the deployed signal is a levelized cost; a case where the dual is
  informative --- or an explicit average-cost interpretation carried through the
  theory --- closes the gap between the deployed price and the marginal-cost ideal.
\item \textbf{The unit/deflator bridge.} Auditing the currency translation, channel by
  channel, is what turns carbon and investment levels from illustrative into
  quotable.
\item \textbf{Loop closure.} The re-solve seam is built and tested; what remains is
  the outer iteration controller --- damping and convergence on the price/quantity
  fixed point --- after which the energy plan reacts to the economy it shapes.
\item \textbf{A calibrated health translation.} The country dose--response and a
  morbidity magnitude replace the placeholder, on the sourced Global Burden of
  Disease scaffolding already in place.
\item \textbf{Energy as a priced production input.} The structural resolution of the
  transmission ambiguity documented in \Cref{sec:discussion} --- the single change
  that would let the model, rather than an assumption, decide how system costs reach
  the economy.
\item \textbf{Deployment behind the planning interface.} The link runs from its own
  environment against the energy model's own solver and case store, so the natural
  end state is as the economic engine behind the planning tool's interface, invoked
  after each energy solve.
\end{enumerate}
